% paper.pdf - the one fixture that is real pdflatex output. Build with
% make_paper.sh (pdflatex twice, so the table and equation references resolve).
%
% What it exercises: LaTeX display equations (amsmath), a booktabs table with
% no vertical rules and only three horizontal ones, a numbered heading
% hierarchy, a footnote set in LaTeX's run-on form ("1We thank" with a raised
% label), an itemize list, and a thebibliography block. Keep it to two pages
% and keep the packages to the ones a base TeX Live install carries.
\documentclass[11pt]{article}
\usepackage[T1]{fontenc}
\usepackage[utf8]{inputenc}
\usepackage{amsmath}
\usepackage{booktabs}
\usepackage[margin=1in]{geometry}
\pagestyle{plain}

\title{A Weight-Free Stand-In for Learned Layout Analysis}
\author{Ada Lovelace \and Charles Babbage}
\date{January 2026}

\begin{document}
\maketitle

\begin{abstract}
We describe a converter that reproduces the heuristic half of a learned
document-analysis pipeline on the PDF text layer alone. It needs no model
weights and runs on a CPU. This fixture is real \LaTeX{} output and exists so
that display equations, a booktabs table and LaTeX-style footnotes are covered
by the regression set.
\end{abstract}

\section{Introduction}

The character stream of a well-formed PDF already encodes reading order, and
on text-layer pages a learned reading-order head rarely improves on it. We
therefore keep stream order and spend the effort on the places where the
stream is silent: headings, tables, footnotes and equations.\footnote{We thank
the reviewers for pressing on this point; the argument is developed in
Section~\ref{sec:method}.}

\subsection{Scope}

The scope is deliberately narrow: journal articles with a text layer,
converted for reading by a language model. Books, slides and forms have
layouts these heuristics were never shown.

\section{Method}
\label{sec:method}

Let $y_i$ denote the observed response for document $i$ and $x_i$ the number
of layout features it contains. We fit
\begin{equation}
  y_i = \alpha + \beta x_i + \varepsilon_i, \qquad
  \varepsilon_i \sim \mathcal{N}(0, \sigma^2),
  \label{eq:model}
\end{equation}
and report the pooled estimate
\begin{equation}
  \hat{\mu} = \frac{1}{n} \sum_{i=1}^{n} y_i .
  \label{eq:mean}
\end{equation}
Equation~\eqref{eq:model} is linear in the parameters; Equation~\eqref{eq:mean}
is what Table~\ref{tab:results} reports in its second column.

The processors run in a fixed order:
\begin{itemize}
  \item line-number stripping, then paragraph reflow;
  \item footnote relabelling before marginalia, so the footer rule cannot
        swallow a note;
  \item heading levels from section numbers, else from clustered line heights.
\end{itemize}

\section{Results}

\begin{table}[h]
  \centering
  \caption{Conversion accuracy by layout feature. $N$ is the number of
           documents in which the feature occurs.}
  \label{tab:results}
  \begin{tabular}{lccr}
    \toprule
    Feature       & $\hat{\mu}$ & Recall & $N$ \\
    \midrule
    Two columns   & 0.98 & 0.97 & 12 \\
    Running head  & 1.00 & 0.94 & 12 \\
    Footnotes     & 0.91 & 0.88 &  9 \\
    Booktabs table& 0.87 & 0.83 &  6 \\
    Display math  & 0.99 & 0.99 & 12 \\
    \bottomrule
  \end{tabular}
\end{table}

Every threshold in the code was set by looking at a failure, not by tuning
against a corpus. That makes the thresholds easy to defend and easy to revise
when a new failure appears. Multi-line table cells remain the weakest point
and are listed as partial support.

\subsection{Limitations}

Inline mathematics such as $\beta > 0$ is not detected; the reference
implementation needs a language model for that as well. Scanned pages inherit
every error of the recogniser.

\begin{thebibliography}{9}
\bibitem{marker} Datalab. Marker: convert PDF to Markdown quickly with high
  accuracy. Software, 2024.
\bibitem{mupdf} Artifex Software. MuPDF and PyMuPDF. Software, 2025.
\end{thebibliography}

\end{document}
